% =====================================================================
% OntoMPO: Ontological Analysis and Refinement of the Model for Project
% Observatories -- SBSI 2027, Trilha de Pesquisa em SI.
%
% Scaffold derivado de template-latex/sbc-template.tex. Compila em branco.
% As nove secoes, seus labels do resumo estruturado e o orcamento de paginas
% estao em artigo/esqueleto.md; os nomes de secao aqui sao os mesmos, na mesma
% ordem.
%
% Diferencas deliberadas em relacao ao template original:
%   - O \usepackage[latin1]{inputenc} duplicado foi REMOVIDO. O template traz
%     utf8 e latin1 nesta ordem, e o segundo vence: com ele, todo acento do
%     corpo em portugues quebra. Nao reintroduzir.
%   - Sem \author, sem \address e sem agradecimento: a revisao e duplamente
%     anonima. O bloco de titulo foi reduzido para nao deixar linha vazia nem
%     marcador de instituicao no lugar deles.
%   - \usepackage[T1]{fontenc} acrescentado. Sem ele o pdflatex compoe cada
%     acento como glifo sobreposto: o PDF sai certo aos olhos, mas a extracao
%     de texto devolve "Introduc,a~o" -- e e' sobre o texto extraido que a
%     varredura de anonimato do ticket 14 trabalha.
%   - A opcao de babel passou de 'brazil', deprecada e ruidosa no log, para
%     'brazilian'.
%   - \nocite{*} e provisorio, so do scaffold: forca todas as entradas de
%     referencias.bib a passarem pelo sbc.bst, de modo que uma entrada quebrada
%     apareca agora e nao na vespera da submissao. O ticket 13 remove esta linha
%     quando as citacoes reais existirem.
% =====================================================================

\documentclass[12pt]{article}

\usepackage{sbc-template}
\usepackage{graphicx,url}
\usepackage[utf8]{inputenc}
\usepackage[T1]{fontenc}
\usepackage[brazilian]{babel}

\sloppy

\title{OntoMPO: Ontological Analysis and Refinement of the\\ Model for Project Observatories}

% Revisao duplamente anonima: sem autoria e sem instituicao.
\author{}
\address{}

\makeatletter
% Bloco de titulo do sbc-template sem as linhas de autor e de instituicao.
\renewcommand{\@maketitle}{\newpage
  \begin{center}
    \vspace*{-.7cm}
    {\XVIPT\bf\@title\par}
    \vglue 6pt plus 3pt minus 3pt
  \end{center}\par}
\makeatother

\begin{document}

\maketitle

% ---------------------------------------------------------------------
% Frontmatter -- orcamento 1,00 pagina. Ticket 13.
% O abstract abaixo e' o texto congelado no JEMS3 em
% submissao/jems3-pacote-registro.md (campo: resumo). Depois de 14/09/2026 o
% registro nao pode mais ser alterado: qualquer divergencia entre este texto e
% o registrado e' achado de conformidade do ticket 14. Nao editar sem editar la'.
% ---------------------------------------------------------------------
\begin{abstract}
  \textbf{Research Context.} Project observatories are organizational units,
  processes or information systems that continuously observe projects,
  centralizing and disseminating information to sustain transparency and
  accountability. The Model for Project Observatories (MPO), the reference
  conceptual model of this domain, organizes 61 concepts in three levels and is
  described in natural language.

  \textbf{Scientific and/or Practical Problem.} Natural-language reference
  models carry representational deficiencies invisible to those who apply them.
  Two initiatives may declare adherence to MPO while adopting semantically
  incompatible interpretations, so adherence is unverifiable and comparison
  between observatories rests on terminological coincidence.

  \textbf{Proposed Solution and/or Analysis.} MPO is submitted to an
  ontological analysis grounded in the Unified Foundational Ontology; the
  resulting revised model is delivered as the artifact. The contribution is
  that analysis and the revised model, not the formalization of MPO as
  published.

  \textbf{Related IS Theory.} Representation Theory (Wand \& Weber) supplies
  the typology used to classify every finding: construct overload, redundancy,
  excess and deficit.

  \textbf{Research Method.} Design Science Research frames the work, continuing
  the cycle that produced MPO; Methontology guides artifact construction. The
  revised OntoUML model is transformed into OWL through gUFO, instantiated over
  an observatory scenario from the literature and interrogated by seven domain
  competency questions in SPARQL, with syntactic and pitfall verification
  before and after revision.

  \textbf{Summary of Results.} The analysis exposed nine representational
  deficiencies: construct overload, invalid mereological relations, semantic
  collapse and missing event and agent microtheories, each classified and
  corrected. Every competency question returned non-empty, semantically correct
  answers, and automated verification reports improved between runs.

  \textbf{Contributions and Impact to IS area.} The work delivers a verifiable
  semantic basis for auditing adherence to MPO, provenance traceability for
  accountability, and a reusable procedure for auditing natural-language
  reference models, addressing semantic interoperability among information
  ecosystems, a named challenge of the Brazilian IS agenda.

  \textbf{Keywords.} Project Observatory; Ontological Analysis; Representation
  Theory; Conceptual Modeling; Design Science Research.
\end{abstract}

\begin{resumo}
  % TODO ticket 13: versao em portugues do resumo estruturado, com os mesmos
  % sete labels. O corpo do artigo e' em portugues, entao o template SBC pede
  % abstract e resumo.
\end{resumo}

% ---------------------------------------------------------------------
% 1. Introducao -- Research Context; Scientific and/or Practical Problem.
% Orcamento 1,25 pagina. Ticket 11.
% ---------------------------------------------------------------------
\section{Introdução}
\label{sec:introducao}

% ---------------------------------------------------------------------
% 2. Referencial Teorico -- Research Context; Related IS Theory.
% Orcamento 2,50 paginas. Ticket 11.
% ---------------------------------------------------------------------
\section{Referencial Teórico}
\label{sec:referencial}

\subsection{Observatórios de Projetos}
\label{sec:observatorios}

\subsection{O Modelo para Observatórios de Projetos}
\label{sec:mpo}

\subsection{Representation Theory}
\label{sec:representation-theory}

\subsection{UFO e OntoUML}
\label{sec:ufo}

% ---------------------------------------------------------------------
% 3. Trabalhos Relacionados -- Research Context.
% Orcamento 0,75 pagina. Ticket 11.
% ---------------------------------------------------------------------
\section{Trabalhos Relacionados}
\label{sec:relacionados}

% ---------------------------------------------------------------------
% 4. Metodo de Pesquisa -- Research Method.
% Orcamento 1,25 pagina. Ticket 11.
% ---------------------------------------------------------------------
\section{Método de Pesquisa}
\label{sec:metodo}

% ---------------------------------------------------------------------
% 5. Analise Ontologica do MPO -- Proposed Solution and/or Analysis.
% Orcamento 3,25 paginas. Ticket 10. Nao e' candidata a corte de paginas:
% e' onde a contribuicao mora.
% ---------------------------------------------------------------------
\section{Análise Ontológica do MPO}
\label{sec:analise}

% ---------------------------------------------------------------------
% 6. A OntoMPO Revisada -- Proposed Solution and/or Analysis.
% Orcamento 1,75 pagina. Ticket 12.
% ---------------------------------------------------------------------
\section{A OntoMPO Revisada}
\label{sec:ontompo}

% ---------------------------------------------------------------------
% 7. Avaliacao -- Research Method; Summary of Results.
% Orcamento 1,75 pagina. Ticket 12.
% ---------------------------------------------------------------------
\section{Avaliação}
\label{sec:avaliacao}

% ---------------------------------------------------------------------
% 8. Discussao e Limitacoes -- Summary of Results; Contributions and Impact
% to IS area. Orcamento 1,00 pagina. Ticket 12.
% ---------------------------------------------------------------------
\section{Discussão e Limitações}
\label{sec:discussao}

% ---------------------------------------------------------------------
% 9. Conclusao e Trabalhos Futuros -- Contributions and Impact to IS area.
% Orcamento 0,75 pagina. Ticket 12.
% ---------------------------------------------------------------------
\section{Conclusão e Trabalhos Futuros}
\label{sec:conclusao}

% ---------------------------------------------------------------------
% Declaracao de uso de IA generativa -- exigida pelo Codigo de Conduta da SBC.
% Orcamento 0,25 pagina. Ticket 12. Sem numeracao e sem identificar autoria.
% ---------------------------------------------------------------------
\section*{Declaração de Uso de IA Generativa}
\label{sec:declaracao-ia}

% ---------------------------------------------------------------------
% Referencias -- orcamento 1,50 pagina.
% \nocite{*} e' provisorio: ver o cabecalho deste arquivo. Ticket 13 remove.
% ---------------------------------------------------------------------
\nocite{*}
\bibliographystyle{sbc}
\bibliography{referencias}

% ---------------------------------------------------------------------
% Apendice A -- consultas SPARQL e resultados completos.
% Orcamento 1,00 pagina. Ticket 13.
% ---------------------------------------------------------------------
\section*{Apêndice A -- Consultas SPARQL e Resultados}
\label{sec:apendice-sparql}

\end{document}
